\section{Milan, conversion, baptism, and return}\label{sec:aug-conversion}

\subsection{Ambrose and a newly intelligible Scripture}

Augustine first approached Ambrose as an accomplished speaker. He found
something more consequential than technique. Ambrose's preaching showed how
Scripture could be read spiritually without making the Old Testament a crude
account of God in bodily form. Augustine did not become Ambrose's intimate
disciple; the \emph{Confessions} repeatedly notes the bishop's inaccessibility
amid pastoral business. The relationship was nonetheless decisive. Ambrose
embodied an intellectually serious Catholic Christianity and a bishop capable
of confronting imperial power.

Milan also intensified Augustine's social ambitions. Monnica arrived and
worked toward a suitable marriage. Because the prospective bride was too
young, the plan required delay. Augustine's long-term partner returned to
Africa after vowing continence, while Adeodatus remained with his father
(\emph{Confessions} 6.13--15). Augustine describes himself as wounded yet
unable to embrace continence, taking another temporary partner. The scene
deserves moral and social clarity. The unnamed woman bears the human cost of a
career marriage project organized by others; the surviving narrative is
Augustine's, not hers. Monnica's role in the plan also resists sentimental
portraiture. Her faith was persistent, but her historical choices were those
of a late Roman mother seeking a son's advancement and Christian stability.

Augustine moved through Academic skepticism toward books he calls
“Platonist,” probably Latin translations or adaptations of Plotinus and
Porphyry. In \emph{Confessions} 7 they gave him conceptual resources for
thinking of God as immaterial and evil as a privation rather than an
independent substance. They did not, in his later judgment, give the humility
of the Incarnation or the grace to live what intellect could glimpse. Paul,
especially Romans, then became central. This intellectual conversion cannot be
separated from the moral crisis narrated in the next book, but neither should
it be compressed into one instant.

\subsection{The garden as remembered prayer}

In late summer 386 Augustine visited the Christian elder Simplicianus and
heard conversion histories involving Victorinus. The courtier Ponticianus
later described Antony and ascetic conversions (\emph{Confessions} 8.2--6).
These narratives made a form of life imaginable and exposed Augustine's
division. In a garden at his lodging, with Alypius nearby, he wept under a fig
tree. He reports hearing a child's sing-song voice repeat \emph{tolle lege}.
Taking the words as a direction to pick up a book and read, he opened Paul's letters and
read Romans 13:13--14; Alypius joined the decision
(\emph{Confessions} 8.12.28--30).

The responsible historical sentence is “Augustine reports,” not “a child
certainly sang for this purpose.” The scene is written as providential
memory, with echoes of Antony's conversion and biblical call narratives.
Literary shaping does not make the event fictitious; it explains how Augustine
wants the reader to perceive it. Earlier Cassiciacum dialogues confirm that a
major redirection had occurred before the \emph{Confessions} was composed.
Modern attempts to identify an exact calendar day, often 15 August 386, go
beyond Augustine's text. “Late summer 386” preserves the secure sequence.

Augustine resigned his rhetorical post after the vacation period, explaining
that illness and his new purpose converged. At the country estate called
Cassiciacum, whose exact modern location remains disputed, a group including
Monnica, Adeodatus, Alypius, relatives, pupils, and friends pursued philosophy,
prayer, and common life. Dialogues such as \emph{Against the Academics},
\emph{On the Happy Life}, \emph{On Order}, and the \emph{Soliloquies} show a
newly committed Christian still working with classical forms. Conversion did
not switch off the intellectual past; it reordered it.

\subsection{Baptism, Ostia, and loss}

Augustine, Adeodatus, and Alypius enrolled for baptism under Ambrose. They were
baptized at the Easter Vigil of 387 (\emph{Confessions} 9.6.14). The checked
source does not supply a modern civil-day conversion, so this revision does
not add one. The later story that Ambrose and Augustine spontaneously composed
the \emph{Te Deum} at this baptism has no contemporary support and belongs to
the tradition review.

The company began a return to Africa. At Ostia, Augustine and Monnica shared
the contemplative conversation remembered in \emph{Confessions} 9.10.23--26,
ascending in speech from created things toward the wisdom beyond them. Soon
afterward Monnica fell ill and died, about fifty-six years old
(\emph{Confessions} 9.11--13). Augustine reports both disciplined restraint and
eventual tears. The episode is neither a modern stenogram nor disposable
decoration. It is the literary and theological farewell through which the son
who had long been the object of his mother's prayer gives her a place within
the Church's hope.

Political and travel conditions delayed the African return. Augustine spent a
period in Rome and wrote against Manichaeans before reaching Thagaste, probably
in 388. There he sold or otherwise disposed of family property and formed a
lay ascetic community with friends. “Monastery” is appropriate if it does not
conjure a later medieval abbey. The community joined prayer, study, material
sharing, and preparation for service. Adeodatus died during this early African
period. Augustine's desired life was now celibate and common, but it would not
remain withdrawn.

\begin{studyframe}
\textbf{Conversion in stages.} The \emph{Confessions} distinguishes a long
intellectual reorientation, the garden decision, catechumenal preparation,
baptism, bereavement, and the construction of a common life. Treating only the
garden as “the conversion” makes a vivid scene replace the historical
process Augustine himself narrates.
\end{studyframe}
